\chapter{Time Synchronisation}\label{ch:nw:time-synchronisation}

European rules require a firm that trades with a high-frequency algorithmic technique to keep the clocks that timestamp its
orders within 100 microseconds of UTC, and to timestamp to the microsecond. The \omterm{def:nw:time-synchronisation:protocols}{Network Time Protocol}'s specification assumes that
an ordinary computer clock may run off frequency by 15 parts per million: fifteen microseconds gained or lost every second. Such a
clock, set right and then left alone, is out of the rule within seven seconds. Staying inside it is a control loop that runs all
day, and proving that it did is a record that must be kept.

The book has used time everywhere: the exchange and receive timestamps of One Quant Book~1, the hardware timestamps of One Quant
Book~13, the captures of chapter~5 to come. This chapter says where a trading firm's time comes from, how it is carried to every
server across a network that delays it by variable amounts, and what the regulators ask a firm to show.

\section{Clocks, oscillators and drift}

A clock counts the cycles of an oscillator. If the oscillator runs slightly fast, the clock gains; if its rate itself wanders
(temperature, ageing), the gain wanders too.

\begin{definition}[Clock offset, fractional frequency offset]\label{def:nw:time-synchronisation:offset}
The \emph{clock offset}\index{clock offset} $\vartheta(t)$ of a clock is the difference between the time it shows and the reference
time at the same instant. Its \emph{fractional frequency offset}\index{fractional frequency offset} $y_{\mathrm f}(t) =
d\vartheta/dt$ is the rate at which the offset grows, a dimensionless number quoted in parts per million (ppm) or per billion
(ppb): a clock with $y_{\mathrm f} = 1$~ppm gains a microsecond every second.
\end{definition}

\begin{definition}[Allan deviation]\label{def:nw:time-synchronisation:adev}
The \emph{Allan deviation}\index{Allan deviation} $\sigma_y(\tau)$ of a clock is the root mean square of the change in its average
fractional frequency between successive intervals of length $\tau$, divided by $\sqrt2$:
$\sigma_y^2(\tau) = \tfrac12 \E\bigl[(\bar y_{k+1} - \bar y_k)^2\bigr]$. From phase samples $x_i$ spaced $\tau$ apart it is
$\sigma_y^2(\tau) = \E\bigl[(x_{i+2} - 2x_{i+1} + x_i)^2\bigr] / (2\tau^2)$.
\end{definition}

The \omterm{def:nw:time-synchronisation:adev}{Allan deviation} answers the question a firm actually has: over an interval of $\tau$, by how much can the clock's rate be
trusted? Unlike the ordinary standard deviation, it converges for the random-walk wanderings of real oscillators, and its slope
against $\tau$ on a log-log plot identifies the noise: $-1$ for white phase noise (a clock whose readings are jittered but whose rate is
right), $+\tfrac12$ for a random walk of frequency (a rate that wanders). \Cref{fig:nw:time-synchronisation:adev} shows both.

\begin{definition}[Holdover]\label{def:nw:time-synchronisation:holdover}
\emph{Holdover}\index{holdover} is the operation of a clock that has lost its reference and keeps time on its own oscillator from
the last frequency correction it had; the holdover time is how long it stays within a tolerance.
\end{definition}

\begin{proposition}[Holdover time]\label{prop:nw:time-synchronisation:holdover}
A clock entering \omterm{def:nw:time-synchronisation:holdover}{holdover} with offset $\vartheta_0$, residual frequency error $y_0$ and a constant ageing $D$ (frequency change per
unit time) has offset $\vartheta(t) = \vartheta_0 + y_0 t + \tfrac12 D t^2$. It stays within a limit $L > |\vartheta_0|$ until the
first positive root of $|\vartheta(t)| = L$; with no ageing, $t = (L - \vartheta_0)/y_0$ for $y_0 > 0$.
\end{proposition}

\begin{proof}
Integrate $d\vartheta/dt = y_0 + Dt$ from $\vartheta(0) = \vartheta_0$.
\end{proof}

\section{Two-way time transfer: NTP and PTP}

A server learns the time from a master by exchanging messages with it. The difficulty is that a message's arrival time mixes the
clocks' offset with the network's delay; two messages in opposite directions separate them, under one assumption.

\begin{definition}[Network Time Protocol, Precision Time Protocol]\label{def:nw:time-synchronisation:protocols}
The \emph{Network Time Protocol}\index{Network Time Protocol} (NTP) synchronises a client to servers by request-response exchanges
over UDP, typically every few seconds to minutes, with timestamps taken in software. The \emph{Precision Time
Protocol}\index{Precision Time Protocol} (PTP, IEEE~1588) synchronises clocks on a network to a master by frequent exchanges, many a
second, designed for timestamps taken in the network hardware and for switches that take part in the protocol.
\end{definition}

\begin{omfigure}
\begin{tikzpicture}[font=\footnotesize,>=Stealth]
\draw[thick] (0,0) -- (0,-4.4) node[below] {master};
\draw[thick] (5,0) -- (5,-4.4) node[below] {slave};
\draw[->,omDef,thick] (0,-0.5) node[left]{$t_1$} -- node[above,sloped]{Sync (and Follow\_Up with $t_1$)} (5,-1.3) node[right]{$t_2$};
\draw[->,omThm,thick] (5,-2.4) node[right]{$t_3$} -- node[above,sloped]{Delay\_Req} (0,-3.3) node[left]{$t_4$};
\draw[->,gray] (0,-3.5) -- node[below,sloped,font=\scriptsize]{Delay\_Resp (with $t_4$)} (5,-4.2);
\draw[<->,gray] (5.9,-0.5) -- node[right,font=\scriptsize]{$d_{\mathrm{ms}}$ plus offset} (5.9,-1.3);
\draw[<->,gray] (-0.9,-2.4) -- node[left,font=\scriptsize]{$d_{\mathrm{sm}}$ minus offset} (-0.9,-3.3);
\end{tikzpicture}

\omcaption{PTP's delay request-response exchange. $t_1$ and $t_4$ are read on the master's clock, $t_2$ and $t_3$ on the slave's. The
first difference is the path delay from master to slave plus the slave's offset; the second is the path delay back minus the offset.
Time runs downward.}\label{fig:nw:time-synchronisation:exchange}
\end{omfigure}

\begin{proposition}[Two-way offset and its blind spot]\label{prop:nw:time-synchronisation:offset}
Let the master send at $t_1$ (its clock), the slave receive at $t_2$ and reply at $t_3$ (its clock), the master receive at $t_4$,
and let the one-way delays be $d_{\mathrm{ms}}$ and $d_{\mathrm{sm}}$. The estimate
\[
  \hat\vartheta = \tfrac12\bigl[(t_2 - t_1) - (t_4 - t_3)\bigr]
  \quad\text{equals}\quad \vartheta + \tfrac12\,(d_{\mathrm{ms}} - d_{\mathrm{sm}}),
\]
and the round trip $(t_4 - t_1) - (t_3 - t_2) = d_{\mathrm{ms}} + d_{\mathrm{sm}}$ does not depend on $\vartheta$. The two-way method
is exact when the two directions take the same time, and no exchange of messages can measure how much they differ.
\end{proposition}

\begin{proof}
$t_2 = t_1 + d_{\mathrm{ms}} + \vartheta$ and $t_4 = t_3 - \vartheta + d_{\mathrm{sm}}$; substitute. Any $(\vartheta, d_{\mathrm{ms}},
d_{\mathrm{sm}})$ and $(\vartheta + a, d_{\mathrm{ms}} - a, d_{\mathrm{sm}} + a)$ produce the same four timestamps.
\end{proof}

\begin{definition}[Path asymmetry, clock servo]\label{def:nw:time-synchronisation:asym}
The \emph{path asymmetry}\index{path asymmetry} of a timing path is the difference $d_{\mathrm{ms}} - d_{\mathrm{sm}}$ between its
two one-way delays; half of it appears, undetected, in every two-way offset estimate. A \emph{clock servo}\index{clock servo} is the
control loop that turns successive offset estimates into corrections of the slave's phase and frequency.
\end{definition}

The asymmetry is fixed by cables (a longer fibre one way), by transceivers and by switches that treat the two directions
differently; the queueing in each direction adds a varying part, which averaging reduces and the servo filters. A firm finds the fixed
part by calibration against a second, independent reference, and subtracts it by configuration: no protocol can find it.

\omcode{code/firm/clocksync/firm_clocksync.py}{58}{68}{One exchange in \texttt{firm.clocksync}: queueing in both directions, timestamp
noise, and the transparent clocks' correction when there are any.}\label{lst:nw:time-synchronisation:exchange}

\section{Grandmasters, satellite references and holdover}

\begin{definition}[Grandmaster clock, GNSS time reference]\label{def:nw:time-synchronisation:gm}
A \emph{grandmaster clock}\index{grandmaster clock} is the clock at the top of a PTP network, the master of all others. A \emph{GNSS
time reference}\index{GNSS time reference} is a receiver of a global navigation satellite system (GPS, Galileo and others) used as a
source of UTC: each satellite's signal carries the system's time and its offset from UTC as kept by a national laboratory.
\end{definition}

A trading firm's grandmaster is usually a dedicated appliance in each cage: a satellite antenna on the data centre's roof, a cable
down to the receiver, and an oscillator that the receiver disciplines and that keeps time in \omterm{def:nw:time-synchronisation:holdover}{holdover} when the signal is lost
(jamming, a failed antenna, a cut cable). The European rule accepts UTC from a satellite system ``provided that any offset from UTC is
accounted for and removed''. GPS's own performance standard specifies the accuracy of the UTC offset it broadcasts: 30 nanoseconds,
95\% of the time. The rest of the error budget is the firm's: the antenna cable's delay (a fixed, measurable length), the receiver, and
every hop from the grandmaster to the timestamp.

\begin{omfigure}
\begin{tikzpicture}[font=\footnotesize,b/.style={draw,rounded corners=2pt,minimum height=0.65cm,align=center,inner sep=3pt},>=Stealth]
\node[b,fill=gray!10] (sat) at (0,0) {satellites\\ (UTC offset\\ in the signal)};
\node[b,fill=omDef!15] (gm) at (3.9,0) {grandmaster\\ receiver and\\ oscillator};
\node[b,fill=omMeth!18] (bc) at (7.4,0) {boundary clock\\ (switch)};
\node[b,fill=omMeth!10] (tc) at (7.4,-1.9) {transparent clock\\ (switch)};
\node[b] (h1) at (10.8,0.6) {server: card's clock};
\node[b] (h2) at (10.8,-0.6) {server: card's clock};
\node[b] (h3) at (10.8,-1.9) {capture appliance};
\draw[->,dashed] (sat) -- node[above,font=\scriptsize,align=center]{antenna\\ cable} (gm);
\draw[->,thick] (gm) -- node[above,font=\scriptsize]{PTP} (bc);
\draw[->,thick] (bc.east) -- (h1.west);
\draw[->,thick] (bc.east) -- (h2.west);
\draw[->,thick] (bc.south) -- (tc.north);
\draw[->,thick] (tc) -- (h3);
\node[font=\scriptsize,align=left,anchor=north west] at (-1.0,-1.2) {errors along the chain: the broadcast\\ UTC offset, the antenna cable,
  the receiver,\\ each boundary clock, residual asymmetry,\\ the card's clock, the step to the host's clock};
\end{tikzpicture}

\omcaption{A cage's timing chain. The grandmaster takes UTC from satellites and serves PTP; a \omterm{def:nw:time-synchronisation:bc}{boundary clock} is a slave on one side
and a master on the other; a \omterm{def:nw:time-synchronisation:bc}{transparent clock} forwards PTP messages and corrects them for the time they spent inside it. Every link
of the chain adds to the error budget of the timestamps at its end.}\label{fig:nw:time-synchronisation:chain}
\end{omfigure}

\section{Boundary and transparent clocks: removing the network's error}

\begin{definition}[Boundary clock, transparent clock]\label{def:nw:time-synchronisation:bc}
A \emph{boundary clock}\index{boundary clock} is a network device with its own clock that is a PTP slave on one port and a master on
its others: it ends one exchange and starts new ones, so that no PTP message crosses it. A \emph{transparent clock}\index{transparent
clock} forwards PTP messages and adds to a correction field in each the time the message spent inside the device (its residence
time), which the endpoints subtract from their timestamps.
\end{definition}

The queueing in a switch is the largest variable delay on a timing path, and it is not symmetric: the Sync message may wait behind a
market-data burst while the Delay\_Req, going the other way, does not. A \omterm{def:nw:time-synchronisation:bc}{transparent clock} measures that wait and reports it; a
\omterm{def:nw:time-synchronisation:bc}{boundary clock} removes the problem by never letting the message queue across the switch. Both need the switch's own hardware to
timestamp, and both leave the fixed asymmetry of cables and transceivers untouched.

\begin{omfigure}
\begin{tikzpicture}
\begin{axis}[width=11.5cm,height=5.6cm,ymode=log,xmin=0,xmax=600,ymin=3,ymax=200000,
  xlabel={seconds after start},ylabel={largest $|\vartheta|$ in each 5~s (ns, log)},
  ytick={10,100,1000,10000,100000},yticklabels={10,100,1\,000,10\,000,100\,000},
  tick label style={font=\small},label style={font=\small},grid=major,grid style={gray!25},
  legend columns=2,legend style={font=\footnotesize,at={(0.5,-0.3)},anchor=north,draw=none,fill=none,
  /tikz/every even column/.append style={column sep=8pt}},table/col sep=comma]
\addplot[omThm,thick] table[x=t_s,y=sw]{figdata/networks/04-time-synchronisation/offset.csv};
\addplot[omExo,thick] table[x=t_s,y=hw]{figdata/networks/04-time-synchronisation/offset.csv};
\addplot[omMeth,thick] table[x=t_s,y=hw_filter]{figdata/networks/04-time-synchronisation/offset.csv};
\addplot[omDef,very thick] table[x=t_s,y=hw_tc]{figdata/networks/04-time-synchronisation/offset.csv};
\addplot[gray,dashed,domain=0:600] {100000};
\legend{software timestamps,hardware timestamps,hardware and min-delay filter,hardware and transparent clocks}
\end{axis}
\end{tikzpicture}

\omcaption{Simulation: a server's clock (20~ppm off, wandering) disciplined over a path of three switches with a mean
\qty{5}{\micro\second} of queueing per switch and direction, eight exchanges a second. Queueing, not the timestamps, dominates until
\omterm{def:nw:time-synchronisation:bc}{transparent clocks} remove it; the dashed line is the 100-microsecond limit. Model values; data: \texttt{fig\_clock.py} on
\texttt{firm.clocksync} (seeded).}\label{fig:nw:time-synchronisation:offsets}
\end{omfigure}

The tutorial runs the four set-ups of \cref{fig:nw:time-synchronisation:offsets} for an hour. After the first six minutes the largest
offset is \qty{26.3}{\micro\second} with software timestamps and \qty{25.9}{\micro\second} with hardware timestamps: timestamping in
hardware barely helps when the messages queue for microseconds. Keeping, of every eight exchanges, the one with the smallest round
trip (the idea of NTP's clock filter: the least-delayed exchange is the least queued) brings it to \qty{11.1}{\micro\second}.
\omterm{def:nw:time-synchronisation:bc}{Transparent clocks} bring it to \qty{25}{\nano\second}, a thousand times better. With \omterm{def:nw:time-synchronisation:bc}{transparent clocks}, a fixed asymmetry of
\qty{2}{\micro\second} leaves the clock \qty{1}{\micro\second} off, exactly as \cref{prop:nw:time-synchronisation:offset} says, and no
amount of filtering sees it.

\omcode{code/firm/clocksync/firm_clocksync.py}{75}{91}{The simulation: the oscillator wanders, the offset grows, an exchange measures
it, and the servo corrects phase and frequency.}\label{lst:nw:time-synchronisation:loop}

\begin{omfigure}
\begin{tikzpicture}
\begin{axis}[width=11.5cm,height=5.4cm,xmode=log,ymode=log,xmin=0.1,xmax=700,ymin=1e-11,ymax=2e-7,
  xlabel={averaging time $\tau$ (s, log)},ylabel={$\sigma_y(\tau)$ (log)},
  tick label style={font=\small},label style={font=\small},grid=major,grid style={gray!25},
  legend columns=2,legend style={font=\footnotesize,at={(0.5,-0.3)},anchor=north,draw=none,fill=none,
  /tikz/every even column/.append style={column sep=8pt}},table/col sep=comma]
\addplot[omThm,very thick,mark=*,mark size=1.4pt] table[x=tau_s,y=free]{figdata/networks/04-time-synchronisation/adev.csv};
\addplot[omDef,very thick,mark=square*,mark size=1.4pt] table[x=tau_s,y=disciplined]{figdata/networks/04-time-synchronisation/adev.csv};
\legend{free-running oscillator,disciplined (transparent clocks)}
\end{axis}
\end{tikzpicture}

\omcaption{Simulation: \omterm{def:nw:time-synchronisation:adev}{Allan deviation} of the chapter's oscillator left alone and of the same clock disciplined over \omterm{def:nw:time-synchronisation:bc}{transparent
clocks}. The free clock's frequency wanders (slope $+\tfrac12$); the disciplined clock's readings jitter by nanoseconds around a correct
rate (slope $-1$). Below a few seconds (between 2 and 8) the free oscillator is the steadier of the two; above, the servo wins by orders of magnitude.
Model values; data: \texttt{fig\_clock.py}.}\label{fig:nw:time-synchronisation:adev}
\end{omfigure}

\section{The regulatory requirements and proving your clock}

\begin{dated}{2026-09}{What the rules ask of clocks}\label{dat:nw:time-synchronisation:rules}
In the European Union, Commission Delegated Regulation (EU) 2017/574 (``RTS~25'') sets, for members and participants of trading venues
using a high-frequency algorithmic trading technique, a maximum divergence from UTC of 100 microseconds and a timestamp granularity of
1 microsecond or better; for trading venues whose gateway-to-gateway latency is one millisecond or less, the same; for voice and
human-mediated request-for-quote systems, one second. UTC is taken from a timing centre listed by the BIPM or from a satellite system
with its offset removed. Firms must establish \omterm{def:nw:time-synchronisation:trace}{traceability to UTC}, identify the exact point at which each timestamp is applied, and
review compliance at least once a year. In the United States, FINRA Rule 6820 requires industry members reporting to the consolidated
audit trail to synchronise their business clocks within 50 milliseconds of the NIST atomic clock (one second for manual orders),
before each day's open and throughout the day, to log the synchronisation for five years, and to certify compliance.
\end{dated}

\begin{definition}[Traceability to UTC]\label{def:nw:time-synchronisation:trace}
\emph{Traceability to UTC}\index{traceability to UTC} is an unbroken, documented chain of comparisons from a firm's timestamps to UTC,
each with a stated uncertainty, such that the firm can show, for any timestamp, how far from UTC it could have been.
\end{definition}

\begin{method}[Proving your clock]\label{met:nw:time-synchronisation:prove}
\begin{enumerate}
\item Document the chain (\cref{fig:nw:time-synchronisation:chain}): reference, grandmasters, every boundary and \omterm{def:nw:time-synchronisation:bc}{transparent clock}, the
  card's clock of each server, the step to the host's clock and the exact point where each kind of timestamp is taken.
\item Write the error budget link by link, with its evidence (a specification, a calibration, a measurement), and compare its worst
  case with the limit.
\item Monitor continuously: each slave's measured offset and path delay, the grandmaster's lock state and \omterm{def:nw:time-synchronisation:holdover}{holdover}, the card-to-host
  step; alarm well inside the limit.
\item Keep the records for the regulators' retention periods, including every period in \omterm{def:nw:time-synchronisation:holdover}{holdover} or out of tolerance.
\item Calibrate the fixed asymmetries against an independent reference, and re-check after every change of cabling or equipment.
\item Review the whole system at least once a year.
\end{enumerate}
\end{method}

\section{Tutorial: disciplining a clock}\label{tut:nw:time-synchronisation:tut}

\begin{tutorial}
\textbf{Goal.} Discipline a simulated server clock four ways over a loaded network and read what limits each.
\textbf{End state:} \cref{fig:nw:time-synchronisation:offsets,fig:nw:time-synchronisation:adev}, the four largest offsets quoted in the
text, and the asymmetry experiment.
\begin{tutsteps}
\item \textbf{The pieces.} \texttt{firm.clocksync} holds the oscillator, the path (base delay, asymmetry, queueing per switch, timestamp
  noise), one exchange (\cref{lst:nw:time-synchronisation:exchange}) and the servo; \texttt{simulate} runs the loop
  (\cref{lst:nw:time-synchronisation:loop}).
\item \textbf{Four set-ups.} \texttt{nw\_clock.SETUPS} differs only in timestamp noise, filtering and \omterm{def:nw:time-synchronisation:bc}{transparent clocks};
  \texttt{summary(name)} gives each one's largest offset after six minutes.
\item \textbf{Stability.} \texttt{adev()} computes the \omterm{def:nw:time-synchronisation:adev}{Allan deviation} of the free and the disciplined clock.
\item \textbf{The blind spot.} \texttt{asymmetry\_bias(2000)} adds a fixed \qty{2}{\micro\second} asymmetry and returns the offset it
  leaves, against $-a/2$. \texttt{python fig\_clock.py} writes the charts' data.
\end{tutsteps}
\textbf{What to change next.} Raise the queueing to \qty{50}{\micro\second} per switch and find the set-ups that break the
100-microsecond limit; halve the exchange rate and read the transparent-clock offset again.
\end{tutorial}

\section{Build: the clock-synchronisation model}\label{bld:nw:time-synchronisation:clocksync}

\begin{build}
\bfield{Purpose} A model of the firm's timing chain in which the effect of every choice (timestamps, filters, \omterm{def:nw:time-synchronisation:bc}{transparent clocks},
asymmetry, \omterm{def:nw:time-synchronisation:holdover}{holdover}) can be computed before it is bought, and a compliance report of offsets against a rule.
\bfield{Interface} \texttt{firm\_clocksync}: \texttt{Oscillator(freq\_ppm, walk\_ppb, seed)}, \texttt{Path(base\_ns, asym\_ns,
queue\_ns, hops, stamp\_ns, seed)}, \texttt{Servo(kp, ki).update(offset, interval)}, \texttt{exchange}, \texttt{simulate(hours,
interval\_s, osc, path, transparent, kp, ki, seed, min\_filter)}, \texttt{allan\_deviation}, \texttt{holdover\_s}, \texttt{compliance};
C++20 servo twin in \texttt{cpp/}.
\bfield{Rules} Times in nanoseconds, frequencies in ppb; the estimate is the two-way formula and nothing else; \omterm{def:nw:time-synchronisation:bc}{transparent clocks}
correct queueing only; every draw is seeded.
\bfield{Acceptance tests} \texttt{code/firm/clocksync/tests/}: the two-way formula exact without asymmetry and biased by half of it with;
\omterm{def:nw:time-synchronisation:bc}{transparent clocks} remove queueing; the servo converges from a 20~ppm error; the servo fixture reproduced in Python and C++; white phase
noise gives an Allan slope of $-1$; \omterm{def:nw:time-synchronisation:holdover}{holdover} with and without ageing; compliance counting.
\bfield{Stretch} A \omterm{def:nw:time-synchronisation:bc}{boundary clock} with its own servo between master and slave, and the error it adds; the peer-delay mechanism; a
monitoring report in the format a regulator would ask to see.
\end{build}

\begin{omsources}
\item Commission Delegated Regulation (EU) 2017/574 (RTS~25), articles 1 and 4 and annex; FINRA Rule 6820.
\item RFC~5905, \emph{Network Time Protocol Version 4}; linuxptp, \texttt{ptp4l(8)}.
\item U.S. Department of Defense, \emph{GPS Standard Positioning Service Performance Standard}, 5th edition (2020).
\end{omsources}

\section{Exercises}

\begin{exercise}[$\star$]\label{exo:nw:time-synchronisation:1}
A clock runs 15~ppm fast. How long after being set correctly does it breach a 100-microsecond limit, and a 50-millisecond one?
\end{exercise}

\begin{exercise}[$\star$]\label{exo:nw:time-synchronisation:2}
An exchange gives $t_1 = 0$, $t_2 = 5\,300$, $t_3 = 6\,300$, $t_4 = 8\,900$ (nanoseconds). What are the estimated offset and the round
trip?
\end{exercise}

\begin{exercise}[$\star$]\label{exo:nw:time-synchronisation:3}
The fibre from master to slave is 40~metres longer than the fibre back. What offset does the two-way method leave, with a group
index of 1.462?
\end{exercise}

\begin{exercise}[$\star\star$]\label{exo:nw:time-synchronisation:4}
Why did hardware timestamps barely help in \cref{fig:nw:time-synchronisation:offsets}, and why does keeping the least-delayed exchange
of eight help?
\end{exercise}

\begin{exercise}[$\star\star$]\label{exo:nw:time-synchronisation:5}
A grandmaster in \omterm{def:nw:time-synchronisation:holdover}{holdover} has a residual frequency error of 1~ppb and ages by $10^{-5}$~ppb a second. How long does it stay within
\qty{100}{\micro\second}?
\end{exercise}

\begin{exercise}[$\star\star$]\label{exo:nw:time-synchronisation:6}
Read \cref{fig:nw:time-synchronisation:adev}: over which averaging times is the disciplined clock worse than the free one, and why is
that not a problem?
\end{exercise}

\begin{exercise}[$\star\star\star$]\label{exo:nw:time-synchronisation:7}
\emph{Coding.} With \texttt{firm.clocksync}, raise the mean queueing to \qty{50}{\micro\second} per switch and direction and compute the
largest offset after six minutes with hardware timestamps, with the min-delay filter and with \omterm{def:nw:time-synchronisation:bc}{transparent clocks}. Which break the
100-microsecond rule?
\end{exercise}

\begin{exercise}[$\star\star\star$]\label{exo:nw:time-synchronisation:8}
\emph{Find the flaw.} ``Our servers' PTP statistics show offsets below 50~ns all day, so our timestamps are within 50~ns of UTC.''
\end{exercise}

\section{Problem: A Hundred Microseconds to UTC}

\begin{problem}[{Weekend problem --- an error budget and a holdover}]\label{pb:nw:time-synchronisation:1}
A firm's timing chain has these error terms (model values, in nanoseconds): the satellite's broadcast UTC offset 30; the antenna cable's
delay after compensation 20; the grandmaster 40; two \omterm{def:nw:time-synchronisation:bc}{boundary clocks} 50 each; the residual \omterm{def:nw:time-synchronisation:asym}{path asymmetry} after \omterm{def:nw:time-synchronisation:bc}{transparent clocks},
half of \qty{200}{\nano\second}; the card's clock 20; the step from the card's clock to the host's clock 500; the software timestamp
point 1\,500. The rule is 100 microseconds.

\medskip\noindent\textbf{Part I --- The budget.}
\begin{enumerate}
\item What is the worst-case error, summing the terms?
\item What is the root-sum-square, if they are independent?
\item Which two terms dominate, and what share of the worst case are they?
\item What margin does the worst case leave under the rule?
\end{enumerate}

\medskip\noindent\textbf{Part II --- Better timestamps.}
\begin{enumerate}[resume]
\item The firm timestamps orders on the card instead of in software. What are the new worst case and root-sum-square?
\item Why does that also remove the card-to-host step from the budget of those timestamps?
\item Which term dominates now?
\item What would halve the remaining worst case?
\end{enumerate}

\medskip\noindent\textbf{Part III --- \omterm{def:nw:time-synchronisation:holdover}{Holdover}.}
\begin{enumerate}[resume]
\item The antenna fails. The grandmaster's oscillator has a residual frequency error of 1~ppb and ages by $10^{-5}$~ppb a second. How
  long before the grandmaster alone uses up the margin of question~4?
\item A cheaper grandmaster has 50~ppb of residual error and ages by $10^{-3}$~ppb a second. How long?
\item Over a weekend of 60 hours without repair, which one stays compliant?
\item What should the monitoring do during \omterm{def:nw:time-synchronisation:holdover}{holdover}?
\end{enumerate}

\medskip\noindent\textbf{Part IV --- The verdict.}
\begin{enumerate}[resume]
\item State the \emph{named result}: the chain's worst-case error with software and with card timestamps, and the two \omterm{def:nw:time-synchronisation:holdover}{holdover} times.
\item How does the European rule's granularity requirement constrain the software timestamp?
\item Why is a second grandmaster with an independent antenna worth more than a better oscillator?
\item What does traceability require beyond the budget?
\item Which of these terms could an auditor check without trusting the firm's own clock?
\item What would a 20-microsecond asymmetry (a misconfigured link) do to the budget?
\item How does the American rule compare with this budget?
\item In one sentence: what does a firm have to know to say how far its timestamps are from UTC?
\end{enumerate}
\end{problem}

\section{Interview questions}

\begin{interviewq}[$\star$ \iqroles{developer}]\label{iq:nw:time-synchronisation:1}
Derive the two-way offset estimate. What assumption does it make?
\end{interviewq}

\begin{interviewq}[$\star\star$ \iqroles{developer}]\label{iq:nw:time-synchronisation:2}
What is the difference between a \omterm{def:nw:time-synchronisation:bc}{boundary clock} and a \omterm{def:nw:time-synchronisation:bc}{transparent clock}, and why does either matter in a trading network?
\end{interviewq}

\begin{interviewq}[$\star\star$ \iqroles{developer}]\label{iq:nw:time-synchronisation:3}
Why would you choose PTP over NTP for a trading host? When is NTP good enough?
\end{interviewq}

\begin{interviewq}[$\star\star$ \iqroles{developer, researcher}]\label{iq:nw:time-synchronisation:4}
You compare timestamps from two servers and one says a fill arrived before the order that caused it. What do you check?
\end{interviewq}

\begin{interviewq}[$\star\star\star$ \iqroles{developer}]\label{iq:nw:time-synchronisation:5}
How would you prove to a regulator that your timestamps were within 100 microseconds of UTC last Tuesday?
\end{interviewq}

\begin{interviewq}[$\star\star$ \iqroles{developer}]\label{iq:nw:time-synchronisation:6}
The satellite antenna fails on a Friday evening. What happens to your clocks, and what do you do?
\end{interviewq}
